\scanpage{109}
% ms. line 01
\deleted{\unsure{soit|ait}} \inserted{l'on ait} $M(r)\leq e^{(r^\sigma)}$ pour $r$ assez grand.
% Paper watermark, inverted relative to the writing: MADE IN FRANCE (oval); The Best Typewriting Quality.

% ms. line 02
\uline{Théorème 2} (Hadamard-Borel) Si l'ordre $\rho$ de $f(z)$
% ms. line 03
est fini, alors pour tout $\varepsilon>0$, la série
% ms. line 04
\[ \sum{}'\frac{1}{|a_n|^{\rho+\varepsilon}}<+\infty \tag*{(5)} \]
% ms. line 04a
($'$ indique
% ms. line 05
qu'on ne somme pas sur les racines nulles)

% ms. line 06
De plus, on a alors (en supposant $f(0)\ne0$)
% ms. line 07
\[ f(z)=e^{P(z)}\prod_n\left(1-\frac{z}{a_n}\right)e^{\frac{z}{a_n}+\ldots+\frac{z^p}{p a_n^p}} \]
% ms. line 08
où $p\leq\rho$, et $P(z)$ est \deleted{\unsure{un polynôme}} \inserted{un polynôme}
% ms. line 09
de degré $\leq\rho$ \unsure{:|.} en d'autres termes, on
% ms. line 10
a \qquad genre de $f$ $\leq$ ordre de $f$.
% Vertical rule in the left margin alongside the theorem.

% ms. line 11
\uline{Corollaire} \quad Si $\rho<1$, i.e s'il existe un $\unsure{0<\mu|\alpha\mu|\alpha}<1$
% ms. line 12
tel que $M(r)\leq e^{(r^\mu)}$ pour $r$ assez grand,
% ms. line 13
alors $f(z)$ est de genre zéro :
% ms. line 14
\[ f(z)=\mathrm c^{\mathrm{te}}\mathbin{\times}\prod_n\left(1-\frac{z}{a_n}\right)\deleted{\unsure{e^{z/a_n}}} \]

% ms. line 15
\uline{Remarque}. \quad Il est très facile de se convaincre
% ms. line 16
que \qquad ordre de $f$ $\leq$ genre de $f+1$. Donc, si $\omega$ et
% ms. line 17
$g$ sont respectivement le genre et l'ordre, on a
% ms. line 18
\[ \boxed{\omega-1\leq g\leq\omega} \]

% ms. line 19
Comme $\unsure{\gamma|g}$ est entier, il en résulte que $\gamma$
% ms. line 20
est connu dès que $\omega$ l'est, sauf si $\omega$ est
% ms. line 21
entier. En ce cas, on peut avoir $\unsure{p|\gamma|g}=\omega$, ou
% ms. line 21a
$\gamma=\omega-1$.
% archive sheet number: 110
